\chapter{计算实验}
\label{sec:comp_study}

为展示所提出的\algoNotation 框架的实际适用性，我们将其应用于航空时刻表规划——服务网络设计与旅客选择的典型实例。在此场景中，网络节点对应机场，服务资源类型对应飞机机队类型，服务弧对应计划航班，空间航段对应起讫机场对。该应用领域特别适合评估我们的框架，因为它展现了前述章节讨论的所有计算挑战：大量可行行程、通过共享网络容量的市场间强耦合，以及对精确时间分辨率的需求。

我们使用源自边疆航空（F9）的真实世界实例评估\algoNotation 框架的性能，并与使用最先进求解器直接求解综合模型的基线（称为"原始模型"）进行比较。为压力测试可扩展性，我们系统地将规划时域从基准24小时周期的$0.5\times$变化到$3\times$，产生六个规模递增的实例。

\section{实验设置}
边疆航空（F9）案例研究代表了美国高度竞争航空市场中典型的超低成本航空公司（ULCC）网络。

该网络涵盖$60$个机场和$259$条有向航段。该航空公司运营同质的空客A320系列机队，建模为$4$种特定机队子类型以考虑不同的座位容量和运营成本。座位容量主要在$186$至$228$座之间。旅客需求侧包括$382$个起讫点（OD）市场。我们考虑最多$2$次中转和$3$小时最大中转窗口的路径选项。

基准案例（F9-1.0x）以24小时规划时域和$15$分钟时间离散化（每天$96$个区间）构建。选择此分辨率是为了与行业标准的准点性能阈值对齐：在商业航空中，在计划时间$15$分钟内到达被监管机构视为准点。因此，该离散化步长代表了时刻表偏差在运营上有意义的最细时间粒度。该基准案例共生成$475{,}776$条可行行程。将近五十万行程变量与离散机队分配约束耦合产生了极其稠密的数学规划，推动了传统计算资源的极限。

为研究框架如何随问题规模扩展，我们通过变化规划时域构建了五个额外实例。缩放变体采用45分钟离散化步长，以在扩展时域的同时保持时间区间数量可处理。F9-0.5x实例覆盖12小时时域，有$81{,}252$条行程；F9-1.5x扩展到36小时，有$557{,}028$条行程；F9-2.0x到48小时，有$951{,}552$条；F9-2.5x到60小时，有$1{,}032{,}804$条；F9-3.0x到72小时，有$1{,}427{,}328$条行程。所有实例共享相同的$60$个机场、$259$个航段、$4$种机队子类型和$382$个OD市场的网络。

\section{计算性能}
我们比较\algoNotation 算法与原始模型基线的收敛行为和可处理性。每个实例在2小时和12小时计算检查点处进行评估。间隙低于$1\%$视为有效最优；一旦达到此阈值，计时停止。结果总结在表~\ref{tab:f9_horizon_scaling}中。

\begin{table}[htbp]
\centering
\caption{不同规划时域下的性能比较}
\label{tab:f9_horizon_scaling}
\begin{tabular}{lcccc}
\toprule
Case & Raw Gap (2h) & \algoNotation Gap (2h) & Raw Gap (12h) & \algoNotation Gap (12h) \\
\midrule
F9-0.5x & 9.44\%  & 11.79\% & 3.10\%  & 6.11\%  \\
F9-1.0x & 18.40\% & 1.84\%  & 4.75\%  & $<1\%$ (2.1h) \\
F9-1.5x & 13.09\% & $<1\%$ (0.1h) & $<1\%$ (5.7h) & $<1\%$ (0.1h) \\
F9-2.0x & --      & 1.25\%  & 15.26\% & 1.25\%  \\
F9-2.5x & 19.65\% & 1.91\%  & 18.06\% & 1.89\%  \\
F9-3.0x & --      & $<1\%$ (1.3h) & -- & $<1\%$ (1.3h) \\
\bottomrule
\end{tabular}
\vspace{0.5em}
\begin{minipage}{0.88\textwidth}
\footnotesize
注：间隙低于1\%报告为$<1\%$，括号中显示达到阈值的实际时间。"--"表示原始模型在相应时间限制内未返回可用的可行现有解或界。
\end{minipage}
\end{table}

在基准案例（F9-1.0x）上，当具有$475{,}776$个行程变量的完整综合模型直接传递给Gurobi求解器时，根松弛和整数定界过程遭受大规模分数退化。经过连续12小时执行后，原始模型仅将间隙缩小到$4.75\%$。相比之下，\algoNotation 算法通过利用下界模型实现更紧的界并求解一系列紧凑的上界子问题，在仅$2.1$小时内达到有效最优（$<1\%$）。

随着规划时域增长，\algoNotation 的优势变得更加明显。对于两个最大实例（F9-2.0x和F9-3.0x），原始求解器在2小时内未能产生任何可用界——它仍在处理根LP松弛。在12小时时，F9-2.0x和F9-2.5x上的原始模型仍表现出$15\%$--$18\%$的间隙，而F9-3.0x完全未能逃离LP松弛阶段。与此同时，\algoNotation 算法在F9-3.0x上仅$1.3$小时达到$<1\%$，在其余大实例上保持$1\%$--$2\%$的间隙。

值得注意的是，\algoNotation 算法在某些较大实例上的表现\emph{优于}最小实例（F9-0.5x）。F9-1.5x和F9-3.0x案例在$1.5$小时内达到最优，而F9-0.5x在12小时时仍保留$6\%$的间隙。这可归因于时域长度与时间粒度之间的交互：缩放实例中使用的45分钟离散化减少了每个机场的时间点数量，缩小了每次\algoNotation 迭代中求解的松弛模型，使下界更快收敛。F9-0.5x实例仅有16个时间点和12小时时域，为松弛提供了有限的时间灵活性，导致较弱的界。

F9-0.5x实例也是原始模型在两个检查点都优于\algoNotation 的唯一案例。仅有$81{,}252$条行程和紧凑的公式，原始MIP足够小，求解器可以通过分支定界稳步推进。\algoNotation 算法的迭代分解产生的开销在如此小的实例上未能完全摊销。这突出表明\algoNotation 框架专门设计用于——并在——单体公式变得不可处理的大规模实例上最为有益。

F9-1.0x实例的分支定界树未探索节点数绘制在图~\ref{fig:comp_tree}中，清楚地显示原始模型的树持续扩展，而现有解和下界几乎保持不变。
\begin{figure}[htbp]
    \centering
    \includegraphics[width=0.75\textwidth]{figure/img/unexpl_raw_end_both.pdf}
    \caption{模型间分支定界树扩展（未探索节点）随时间的比较。}
    \label{fig:comp_tree}
\end{figure}
相比之下，\algoNotation 算法在分支定界过程的早期阶段完成，树规模小得多。\algoNotation 算法通过利用下界模型（LBM）实现更紧的下界来规避弱LP松弛。这种结构优势随着规划时域增长变得越来越明显：原始模型的树规模爆炸，而\algoNotation 子问题保持紧凑。

\section{优化时刻表分析}
基准案例（F9-1.0x）的全局最优解揭示了调度行为的深刻运营和商业洞察。我们系统地从精确解中提取激活的变量以分析底层运营。

最优时刻表物理上恰好激活$348$个航班，在指定网络中所有$60$个机场保持活跃服务。机队分配根据需求密度完美平衡可用子类型：$178$个计划航班由更高容量的A321-271NX运营，其余$170$个航班分配给A320-214。这种异构部署确保最大飞机严格分配到具有充足合并需求的干线航线，最小化每座英里成本，而较小变体被部署以高效满足次要频率要求而不产生过度容量浪费。

算法不是盲目地试图服务所有潜在OD对，而是表现出敏锐的盈利驱动选择性。在$382$个市场中，时刻表策略性地仅从$120$个最有利可图的市场捕获需求，有意识地放弃结构性不盈利的市场以节约运营成本。在这些目标市场中，基于销售的线性规划（SBLP）组件确保了$5.42\%$的平均市场份额。这高于F9在这些市场的$4.85\%$历史市场份额，表明优化时刻表有效地在最有利可图的细分市场中捕获了更大比例的需求。

此外，优化利用了网络的空间连通性。最优旅客分配积极利用$112$条多段中转行程，而仅有$64$条直达行程。这种对中转行程的高度依赖展示了模型高效拼接来自不同市场的碎片化需求的能力。通过在重点城市间按时间顺序同步物理航班的出发和到达界限，\algoNotation 框架无缝捕获间接需求并最大化全系统运营利润。
